\section{Introduction}

In April 1992, forty-odd people met in Williamsburg, Virginia, because
distributed-memory parallel computing had a portability
problem~\cite{dongarra1993mpi1draft,walker1992williamsburg}. Every vendor
shipped a message-passing library---NX on Intel machines, Vertex on nCUBE,
CMMD on the Connection Machine---and every research group shipped
another---PVM, p4, Chameleon, Zipcode, PARMACS, Express,
TCGMSG, CHIMP~\cite{geist1994pvm,butler1994p4,skjellum1994zipcode,bomans1990parmacs}.
The libraries were not merely different in spelling. They differed in what a
message \emph{was}, in whether a receive could match on a wildcard, in
whether a collective operation existed at all. A parallel application was
tied to a machine, and a parallel \emph{library} was almost impossible to
write, because a library and its caller shared one flat tag space and could
silently steal each other's messages.

The MPI Forum's answer, standardised in MPI-1.0 in May
1994~\cite{mpi-forum-1994,dongarra1996cacm}, was not a better library. It
was a \emph{specification}: a fixed vocabulary of communication operations
with precise semantics, no implementation, and no opinion about what the
ranks computed. Thirty years and five major revisions later, essentially
every distributed-memory scientific code in production speaks
it~\cite{bernholdt2020survey,gropp2001learning}.

\subsection{The same problem, again}

Multi-agent systems built on large language models are in the position
distributed-memory computing occupied in 1992. There are many frameworks,
each with its own model of what an agent is and how agents
communicate. AutoGen rewrote itself around an actor model; LangGraph
compiles a state machine over a shared, checkpointed state object; CrewAI
composes roles and tasks; the OpenAI Agents SDK passes control by handoff;
MetaGPT and ChatDev encode a software-company org
chart~\cite{microsoft2025autogen04,langchaindocs2026persistence,crewaidocs2026,openaiagentsdocs2026handoffs,hong2024metagpt,chatdevissue2023}.
A harness written for one is a rewrite for another, and the reason is
familiar: each of them fuses three separable concerns---what the agents do,
how they are scheduled, and how they exchange information---into one
artifact.

The protocols that have appeared to address agent interoperability do not
close this gap, and mostly do not try to. MCP standardises how one client
reaches one server's tools, resources and
prompts~\cite{mcpblog2026}. A2A standardises how one agent discovers and
delegates a task to another, with agent cards, tasks and
artifacts~\cite{a2adocs2026v1}. ACP and ANP occupy adjacent
positions~\cite{ibm2026acp,ehtesham2025survey}. All of them are, in the vocabulary of
this paper, \emph{point-to-point}: a client, a server, a request, a
response. None of them has a notion of a group, a collective operation, a
barrier, a consistency model for shared state, or a defined outcome when a
participant dies mid-operation. They are the agent world's sockets. What is
missing is its MPI.

That absence has a measurable cost. The largest empirical study of
multi-agent LLM failures to date, which analyses 200+ traces across seven
frameworks, finds that the majority of failures are not reasoning failures
at all but \emph{coordination} failures: agents that disobey a specification
they were given, that lose track of what a peer already decided, that drop a
step because they never learned another agent had finished, that terminate
prematurely, or that never verify an artifact anyone
produced~\cite{cemri2025mast}. Practitioners report the same thing from the
other direction: Cognition's widely-read argument against multi-agent
architectures is essentially that context does not propagate reliably
between agents and that conflicting implicit decisions
compound~\cite{yan2025dontbuild}, while Anthropic's report on a
production multi-agent research system spends most of its engineering
discussion on exactly these problems---coordination, state, and
recovery---rather than on prompting~\cite{anthropic2025multiagent}.

Those are not new problems. They are, respectively, contract violation,
cache coherence, synchronisation, termination detection, and validation. A
field that has been solving them for forty years has names for all of them.

\subsection{What we take from MPI, and what we cannot}

This paper asks what happens if the analogy is taken seriously and pushed as
far as it will go. The answer is: much further than we expected, and then it
breaks in four specific and interesting places.

What transfers is most of MPI's structure, and in particular its single best
idea. A communicator is a group plus an opaque \emph{context}, and the
context guarantees that a message sent on one communicator cannot be
received on another. That property is what made MPI \emph{libraries}
possible: before it, libraries reserved tag ranges by convention, and two
libraries that disagreed produced silent, intermittent
corruption~\cite{skjellum1994zipcode,gropp2001learning}. Multi-agent systems
are in the pre-context era today---the dominant designs are a shared
conversation buffer every participant reads and writes, or handoffs
addressed by agent name---and they exhibit exactly the failure the Forum
identified in 1993.

What does not transfer is the assumption that a rank is a cheap,
deterministic, fail-stop process with reusable memory. Concretely:

\begin{enumerate}
\item \textbf{Context is consumed, not occupied.} An MPI rank may receive a
terabyte over a job's life while never holding more than a megabyte; memory
pressure is a function of the peak live set. Every token an agent reads is
retained for the rest of the conversation, so pressure is a function of
\emph{cumulative ingest}. The right systems analogue is not \code{malloc}
but a quota with admission control.

\item \textbf{Reduction operators are semantic.} \mpiop{SUM} is
associative, commutative, and output-size-preserving. ``Merge these two
drafts'' is none of the three, and getting it wrong does not produce a
rounding difference, it produces a different document.

\item \textbf{Failure is not fail-stop.} Agents crash, but they also stall
while looking perfectly alive, return confidently malformed output, and run
out of money. Conflating these produces the two classic pathologies of agent
harnesses: killing a slow-but-working agent, and waiting forever for a
fast-but-broken one.

\item \textbf{The cost model inverts.} In HPC a message costs microseconds
and a floating-point operation costs sub-nanoseconds, so one optimises
communication. An agent turn costs tens of seconds and real money, while a
message is a file write. Communication is nearly free and computation is
nearly everything.
\end{enumerate}

Each of these forces a specific, minimal change to the MPI design rather
than a wholesale redesign, and \Cref{sec:mismatch,sec:design} develop them
in turn.

\subsection{Contributions}

\begin{itemize}
\item \textbf{A protocol.} \ampi{} (\Cref{sec:design}): communicators with
isolated contexts, groups, token-bounded contract-carrying datatypes, four
send modes including one that completes on \emph{ingestion}, the full
collective catalogue with declared operator algebra, one-sided windows with
reference-by-default reads and a bounded retrieval read, collective
artifact I/O, a failure lattice with revoke/shrink/agree/replace, a
name-shifted profiling interface and a control/performance variable
interface.

\item \textbf{A cost model} (\Cref{sec:cost}) in which the dominant term is
agent turns on the critical path and feasibility is a hard constraint, and
the demonstration that this inverts several of MPI's standard algorithm
choices.

\item \textbf{Two results about operator algebra} that we did not
anticipate. A capacity-aware reduction tree must be planned against the
root's \emph{cumulative} ingest, not its per-round ingest, because an agent
cannot free a buffer (\Cref{sec:cost}); and a prefix scan over shared naming
decisions requires a \emph{non-commutative} merge operator, because a
commutative one cannot express precedence and therefore leaves participants
disagreeing about what they agreed on (\Cref{sec:ops}).

\item \textbf{An implementation} (\Cref{sec:impl}) with a swappable device
layer and a command-line binding, so that any process able to run a shell
command is a rank---which is what makes the protocol usable by agents that
have no library bindings and never will.

\item \textbf{An evaluation with genuine multi-agent runs}
(\Cref{sec:eval}): a twelve-rank book translation, an eight-module
collaborative software build judged by a frozen integration suite, and
fault-injection runs, together with microbenchmarks to 128 ranks.
\end{itemize}

We also report, in \Cref{sec:discussion}, a coordination failure we caused
ourselves: a launcher concurrency limit interacting with a collective's
all-participants requirement produced a deadlock that no rank could observe
locally. It is a good advertisement for the protocol's own diagnostics, and
a reminder that the hard part of a multi-agent system is not the agents.
